\documentclass[../../main.tex]{subfiles}

\graphicspath{{../../image/}{../../figure/}}

\begin{document}

\section{密度变换}\label{section:ch02-密度变换}

当我们对随机变量进行变换时，概率密度也会相应地发生变化。假设我们有一个随机变量$x$，其概率密度为$p_x(x)$。现在，我们通过一个函数$y = f(x)$对$x$进行变换，得到新的随机变量$y$。我们想知道$y$的概率密度$p_y(y)$是什么。

对于离散变量，变换是直接的：$y$取某个值的概率等于$x$取所有能映射到该值的值的概率之和。

对于连续变量，情况稍微复杂一些。概率密度的变换需要满足：概率质量在变换前后保持不变。也就是说，落在$x$的某个区间内的概率，应该等于落在变换后对应的$y$的区间内的概率。

如果变换$f$是单调的（即可逆的），那么概率密度的变换公式为

\begin{align}
  p_y(y) = p_x(x) \left| \frac{\mathrm{d}x}{\mathrm{d}y} \right|
  \label{align:ch02-密度变换}
\end{align}

其中$x$是$y$的反函数，即$x = f^{-1}(y)$。

绝对值的出现是因为我们关心的是概率质量，而不是有向面积。如果$f$不是单调的，那么我们需要将所有能映射到同一个$y$值的$x$值的贡献加起来。

\end{document}